\chapter{Inverse Matrices and Elementary Matrices}

\section{Inverses}

\begin{definition}
An $n\times n$ matrix $A$ is \emph{invertible} if there is an $n\times n$ matrix $B$ with
$AB=I=BA$. Then $B$ is \emph{the inverse} of $A$, written $A^{-1}$. A non-invertible square matrix is
called \emph{singular}.
\end{definition}

The inverse is unique: if $B$ and $C$ are both inverses then $B=BI=B(AC)=(BA)C=IC=C$.
Only square matrices have (two-sided) inverses.

\begin{theorem}[One-sided suffices]
If $A,B$ are $n\times n$ and $AB=I$, then $BA=I$; so $A$ and $B$ are inverses of each other.
\end{theorem}
\begin{proof}
If $B\vec x=\vec0$ then $\vec x=I\vec x=A(B\vec x)=\vec 0$, so $\Null(B)=\{\vec0\}$ and
$\rank B=n$. Hence $B\vec x=\vec y$ is solvable for every $\vec y$. For any $\vec y$ write
$\vec y=B\vec x$; then $(BA)\vec y=BAB\vec x=B(AB)\vec x=B\vec x=\vec y$. So $BA\vec y=I\vec y$ for all
$\vec y$, and $BA=I$ by the Matrices Equal Theorem.
\end{proof}

\begin{theorem}[Properties]
If $A,B$ are invertible $n\times n$ and $c\neq 0$:
\begin{enumerate}
  \item $(A^{-1})^{-1}=A$ and $(cA)^{-1}=\frac1cA^{-1}$;
  \item $(AB)^{-1}=B^{-1}A^{-1}$ (note the order);
  \item $(A^T)^{-1}=(A^{-1})^T$;
  \item $(A^k)^{-1}=(A^{-1})^k$.
\end{enumerate}
\end{theorem}
\begin{proof}[Proof of (2) and (3)]
$(AB)(B^{-1}A^{-1})=A(BB^{-1})A^{-1}=AA^{-1}=I$. And $A^T(A^{-1})^T=(A^{-1}A)^T=I^T=I$.
\end{proof}

\begin{theorem}[$2\times2$ inverse]
$A=\begin{bmatrix}a&b\\c&d\end{bmatrix}$ is invertible iff $ad-bc\ne0$, and then
\[
  A^{-1}=\frac{1}{ad-bc}\begin{bmatrix}d&-b\\-c&a\end{bmatrix}.
\]
\end{theorem}
\begin{proof}
Multiply out: $\begin{bmatrix}a&b\\c&d\end{bmatrix}\begin{bmatrix}d&-b\\-c&a\end{bmatrix}=(ad-bc)I$.
If $ad-bc=0$, this product shows a non-zero vector (a column of the second matrix, unless
$A=O$) in $\Null(A)$, so $A$ is not invertible.
\end{proof}

\section{Computing inverses by row reduction}

To find $A^{-1}$ we must solve $A\vec b_j=\vec e_j$ for each column $\vec b_j$ of $B=A^{-1}$. All
$n$ systems have the same coefficient matrix, so we solve them at once.

\begin{important}[Matrix Inversion Algorithm]
Row reduce $[A\mid I]$. If the RREF has the form $[I\mid B]$ then $A$ is invertible and
$B=A^{-1}$. If the left block does not reduce to $I$ (i.e.\ $\rank A<n$), then $A$ is not invertible.
\end{important}

\begin{example}
Find the inverse of $A=\begin{bmatrix}1&2&3\\2&5&3\\1&0&8\end{bmatrix}$.
\[
\left[\begin{array}{ccc|ccc} 1&2&3&1&0&0\\ 2&5&3&0&1&0\\ 1&0&8&0&0&1\end{array}\right]
\xrightarrow{\substack{R_2-2R_1\\R_3-R_1}}
\left[\begin{array}{ccc|ccc} 1&2&3&1&0&0\\ 0&1&-3&-2&1&0\\ 0&-2&5&-1&0&1\end{array}\right]
\]
\[
\xrightarrow{R_3+2R_2}
\left[\begin{array}{ccc|ccc} 1&2&3&1&0&0\\ 0&1&-3&-2&1&0\\ 0&0&-1&-5&2&1\end{array}\right]
\xrightarrow{-R_3}
\left[\begin{array}{ccc|ccc} 1&2&3&1&0&0\\ 0&1&-3&-2&1&0\\ 0&0&1&5&-2&-1\end{array}\right]
\]
\[
\xrightarrow{\substack{R_1-3R_3\\R_2+3R_3}}
\left[\begin{array}{ccc|ccc} 1&2&0&-14&6&3\\ 0&1&0&13&-5&-3\\ 0&0&1&5&-2&-1\end{array}\right]
\xrightarrow{R_1-2R_2}
\left[\begin{array}{ccc|ccc} 1&0&0&-40&16&9\\ 0&1&0&13&-5&-3\\ 0&0&1&5&-2&-1\end{array}\right].
\]
So $A^{-1}=\begin{bmatrix}-40&16&9\\13&-5&-3\\5&-2&-1\end{bmatrix}$. Check the first row of $AA^{-1}$:
$(1,2,3)\cdot(-40,13,5)=-40+26+15=1$, $(1,2,3)\cdot(16,-5,-2)=16-10-6=0$,
$(1,2,3)\cdot(9,-3,-1)=9-6-3=0$.
\end{example}

\begin{example}[A singular matrix]
\[
\left[\begin{array}{cc|cc} 1&2&1&0\\ 2&4&0&1\end{array}\right]
\xrightarrow{R_2-2R_1}
\left[\begin{array}{cc|cc} 1&2&1&0\\ 0&0&-2&1\end{array}\right].
\]
The left block has rank $1<2$, so $\begin{bmatrix}1&2\\2&4\end{bmatrix}$ is not invertible (consistent
with $ad-bc=4-4=0$).
\end{example}

If $A$ is invertible, then $A\vec x=\vec b$ has the unique solution $\vec x=A^{-1}\vec b$. (In practice
row reducing $[A\mid\vec b]$ is cheaper than computing $A^{-1}$.)

\section{Elementary matrices}

\begin{definition}
An \emph{elementary matrix} is a matrix obtained from $I_n$ by performing exactly one elementary row
operation.
\end{definition}

\begin{theorem}
If $E$ is the elementary matrix for a row operation, then $EA$ is the result of performing that
operation on $A$. Every elementary matrix is invertible, and $E^{-1}$ is the elementary matrix of the
reverse operation.
\end{theorem}

\begin{example}
In $3\times 3$: $E=\begin{bmatrix}1&0&0\\-2&1&0\\0&0&1\end{bmatrix}$ performs $R_2-2R_1$, with
$E^{-1}=\begin{bmatrix}1&0&0\\2&1&0\\0&0&1\end{bmatrix}$ (performing $R_2+2R_1$).
\end{example}

\begin{theorem}
If $A$ row reduces to $R$ via operations with elementary matrices $E_1,\dots,E_k$, then
$E_k\cdots E_2E_1A=R$. In particular if $A$ is invertible then $R=I$, and
\[
  A^{-1}=E_k\cdots E_1, \qquad A = E_1^{-1}E_2^{-1}\cdots E_k^{-1},
\]
so every invertible matrix is a product of elementary matrices.
\end{theorem}

This also explains the inversion algorithm: applying $E_k\cdots E_1=A^{-1}$ to $[A\mid I]$ gives
$[I\mid A^{-1}]$.

\begin{example}
Write $A=\begin{bmatrix}2&4\\1&3\end{bmatrix}$ as a product of elementary matrices.
\[
A\xrightarrow{R_1\leftrightarrow R_2}\begin{bmatrix}1&3\\2&4\end{bmatrix}
\xrightarrow{R_2-2R_1}\begin{bmatrix}1&3\\0&-2\end{bmatrix}
\xrightarrow{-\frac12R_2}\begin{bmatrix}1&3\\0&1\end{bmatrix}
\xrightarrow{R_1-3R_2}\begin{bmatrix}1&0\\0&1\end{bmatrix}.
\]
With $E_1=\begin{bmatrix}0&1\\1&0\end{bmatrix}$, $E_2=\begin{bmatrix}1&0\\-2&1\end{bmatrix}$,
$E_3=\begin{bmatrix}1&0\\0&-\frac12\end{bmatrix}$, $E_4=\begin{bmatrix}1&-3\\0&1\end{bmatrix}$ we get
$E_4E_3E_2E_1A=I$, so
\[
  A=E_1^{-1}E_2^{-1}E_3^{-1}E_4^{-1}
   =\begin{bmatrix}0&1\\1&0\end{bmatrix}\begin{bmatrix}1&0\\2&1\end{bmatrix}
    \begin{bmatrix}1&0\\0&-2\end{bmatrix}\begin{bmatrix}1&3\\0&1\end{bmatrix}.
\]
Check: the last two give $\begin{bmatrix}1&3\\0&-2\end{bmatrix}$; times
$\begin{bmatrix}1&0\\2&1\end{bmatrix}$ on the left gives $\begin{bmatrix}1&3\\2&4\end{bmatrix}$; swapping
rows gives $A$. Also $A^{-1}=E_4E_3E_2E_1=\frac12\begin{bmatrix}3&-4\\-1&2\end{bmatrix}$, matching the
$2\times2$ formula with $ad-bc=2$.
\end{example}

\section{The Invertible Matrix Theorem}

\begin{theorem}[Invertible Matrix Theorem]
For an $n\times n$ matrix $A$, the following are equivalent:
\begin{enumerate}
  \item $A$ is invertible;
  \item the RREF of $A$ is $I$;
  \item $\rank A=n$;
  \item $A\vec x=\vec b$ is consistent with a unique solution for every $\vec b\in\R^n$;
  \item $A\vec x=\vec 0$ has only the trivial solution, i.e.\ $\Null(A)=\{\vec 0\}$;
  \item the columns of $A$ are linearly independent;
  \item the columns of $A$ span $\R^n$ (form a basis of $\R^n$);
  \item the rows of $A$ form a basis of $\R^n$;
  \item $A^T$ is invertible;
  \item $L(\vec x)=A\vec x$ is one-to-one; equivalently, onto;
  \item $A$ is a product of elementary matrices;
  \item $\det A\neq 0$ (Chapter 8);
  \item $0$ is not an eigenvalue of $A$ (Chapter 9).
\end{enumerate}
\end{theorem}
\begin{proof}[Proof sketch]
(2)$\Leftrightarrow$(3) by definition of rank for square matrices. (3)$\Leftrightarrow$(4)$\Leftrightarrow$(5)
by the System-Rank Theorem ($n-\rank A=0$ parameters and $\rank A=n=m$). (5)$\Leftrightarrow$(6) and
(3)$\Leftrightarrow$(7) are the row-reduction criteria for independence and spanning. (1)$\Rightarrow$(5):
$A\vec x=\vec0$ gives $\vec x=A^{-1}\vec0=\vec 0$. (2)$\Rightarrow$(11)$\Rightarrow$(1) by the previous
section. (1)$\Leftrightarrow$(9) from $(A^T)^{-1}=(A^{-1})^T$, and (8) is (6) for $A^T$.
\end{proof}

\begin{important}[Pitfalls]
$(A+B)^{-1}\neq A^{-1}+B^{-1}$ in general (e.g.\ $A=B=I$ gives $(2I)^{-1}=\frac12I\ne 2I$).
And do not ``divide'' by matrices: from $AB=AC$ you may conclude $B=C$ \emph{only} if $A$ is
invertible (multiply on the left by $A^{-1}$).
\end{important}

\section*{Summary}
\begin{itemize}
  \item $A^{-1}$ satisfies $AA^{-1}=A^{-1}A=I$; one side suffices for square matrices; it is unique.
  \item $(AB)^{-1}=B^{-1}A^{-1}$, $(A^T)^{-1}=(A^{-1})^T$; $2\times 2$ formula with $ad-bc$.
  \item Row reduce $[A\mid I]$ to $[I\mid A^{-1}]$.
  \item Elementary matrices perform row operations by left multiplication; invertible matrices are
        products of elementary matrices.
  \item Invertible Matrix Theorem: invertible $\Leftrightarrow$ rank $n$ $\Leftrightarrow$ trivial
        null space $\Leftrightarrow$ columns form a basis $\Leftrightarrow\det A\ne0$.
\end{itemize}
